%% TRANSCRIBED (PDF 111-121 = printed III-20..III-30)
%% NOTE(p121): the tail of §3.4 (end of the verification of (b), the "End
%% of the proof of the theorem" block and the Grössencharakter Exercise)
%% is printed at the top of III-30 (PDF 121), BEFORE the "APPENDIX" head;
%% it belongs to this §3 and is transcribed at the end of this file (the
%% Appendix itself starts in ch3_app.tex).
%% NOTE(p115): the first item of Proposition 1 is printed "(1.)" (with a
%% stray period) while items 2 and 3 are "(2)", "(3)"; transcribed as
%% printed.

\sechead{3}{The case of a composite of quadratic fields}%
  {THE CASE OF A COMPOSITE OF QUADRATIC FIELDS}

\subsechead{3.1}{Statement of the result}

The aim of this \S{} is to prove:

\begin{theorem}
Let $\rho$ be a rational, semi-simple, \textit{l}-adic abelian
representation of $K$.  Assume
\[
(*)\quad K\text{ is a composite of quadratic extensions of }\Q.
\]
Then $\rho$ is locally algebraic (and hence stems from a linear
representation of some $S_m$, cf.~\S\ref{sec:III.2.3}).
\end{theorem}

This applies in particular when $K=\Q$ or when $K$ is quadratic over
$\Q$.

\begin{remarks}
\begin{nlist}
\item An analogous result holds for $E$-rational semi-simple abelian
$\lambda$-adic representations (cf.~\S\ref{sec:III.2.4}).
\item It is quite likely that condition $(*)$ is not necessary.  But
proving this seems to require stronger results on transcendental
numbers than the ones now available.
\end{nlist}
\end{remarks}

\subsechead{3.2}{A criterion for local algebraicity}

\begin{proposition}
Let $\rho\colon \Gal(\bar K/K)^{\mathrm{ab}}\to\Aut(V_\ell)$ be a
rational semi-simple \textit{l}-adic abelian representation of $K$.
Assume that there exists an integer $N\geq 1$ such that $\rho^N$ is
locally algebraic.  Then $\rho$ is locally algebraic.
\end{proposition}

\begin{proof}
Since $\rho$ is semi-simple, it can be brought in diagonal form over a
finite extension of $\Q_\ell$, hence gives rise to a family
$\{\psi_1,\ldots,\psi_n\}$ of $n$ continuous characters
$\psi_i\colon C_K\to\bar\Q_\ell^{\,*}$, where $C_K$ is the
idèle-class group of $K$, and $n=\dim.\,V_\ell$.  Let
$\chi_1=\psi_1^N,\ldots,\chi_n=\psi_n^N$ be the corresponding
characters occurring in $\rho^N$.  Since $\rho^N$ is locally
algebraic, to each $\chi_i$ corresponds an algebraic character
$\chi_i^{\mathrm{alg}}\in X(T)$ of the torus $T$ (here we identify
$X(T)$ with $\Hom(T/\bar\Q_\ell,\Gm/\bar\Q_\ell)$, since
$\bar\Q_\ell$ is algebraically closed).  Each $\chi_i^{\mathrm{alg}}$
is of the form $\prod\limits_{\sigma\in\Gamma}
[\sigma]^{n_\sigma(i)}$, where $\Gamma$ is the set of embeddings of
$K$ into $\bar\Q_\ell$, cf.~Chap.~\ref{chap:II},
\S\ref{sec:II.1.1}.  One has
\[
\chi_i(x)=\chi_i^{\mathrm{alg}}(x^{-1})
=\prod\sigma(x)^{-n_\sigma(i)}
\]
for all $x\in K_\ell^*$ close enough to $1$.
\end{proof}

\begin{lemma}
All the integers $n_\sigma(i)$, $1\leq i\leq n$, $\sigma\in\Gamma$,
are divisible by $N$.
\end{lemma}

\begin{proof}[Proof of the lemma]
Let $U$ be an open subgroup of $\bar\Q_\ell^{\,*}$ containing no
$N^{\mathrm{th}}$-root of unity except $1$, and let $m$ be a modulus
of $K$ such that $\psi_i(x)\in U$ for all $x\in U_m$ and
$i=1,\ldots,n$; the existence of such an $m$ follows from the
continuity of $\psi_1,\ldots,\psi_n$.  We take $m$ large enough so
that:
\begin{alist}[label=\alph*)]
\item It is a modulus of definition for $\rho^N$.
\item $\rho$ is unramified at all $v\in\Supp(m)$, and the
corresponding Frobenius elements $F_{v,\rho}$ have a characteristic
polynomial with rational coefficients.
\end{alist}
Let $K_m$ be the abelian extension of $K$ corresponding to the open
subgroup $K^*U_m$ of the idèle group $I$, and let $L$ be a finite
Galois extension of $\Q$ containing $K_m$.  Choose a prime number $p$
which is distinct from $\ell$, is not divisible by any place of
$\Supp(m)$, and splits completely in $L$.  Let $v$ be a place of $K$
dividing $p$, and let $f_v$ be an idèle which is a uniformizing
element at $v$ and is equal to $1$ elsewhere.  The fact that $v$
splits completely in $K_m$ (since it does in $L$) implies that $f_v$
is the norm of an idèle of $K_m$, hence (by class-field theory)
belongs to $K^*U_m$; this means that the prime ideal
$\mathfrak{p}_v$ is a principal ideal $(\alpha)$, with
$\alpha\equiv 1\bmod.\ m$ and $\alpha$ positive at all real places of
$K$.

Let $x=\psi_i(f_v)$ and $y=\chi_i(f_v)$, so that $y=x^N$; these are
the Frobenius elements of $\psi_i$ and $\chi_i$ relative to $v$.  By
definition of $\chi_i^{\mathrm{alg}}$, we have
\[
y=\chi_i^{\mathrm{alg}}(\alpha)
=\prod_{\sigma\in\Gamma}\sigma(\alpha)^{n_\sigma(i)}
\]
where $\alpha$ is as above.

Hence $y$ belongs to the subfield $\widetilde L$ of $\Q_\ell$
corresponding to $L$ (this field is well defined since $L$ is a
Galois extension of $\Q$).  Moreover, if $w_\sigma$ is any place of
$L$ such that $w_\sigma\circ\sigma$ induces $v$ on $K$, we have (as
in chap.~\ref{chap:II}, \S\ref{sec:II.3.4}):
\[
w_\sigma(y)=n_\sigma(i).
\]

Assume now that $n_\sigma(i)$ is not divisible by $N$.  Then $x$,
which is an $N^{\mathrm{th}}$-root of $y$, does not belong to
$\widetilde L$.  Hence there is a non-trivial
$N^{\mathrm{th}}$-root of unity $z$ such that $x$ and $zx$ are
conjugate over $\widetilde L$, and \emph{a fortiori} over $\Q$.
Since the characteristic polynomial of $F_{v,\rho}$ has rational
coefficients, any conjugate over $\Q$ of an eigenvalue of
$F_{v,\rho}$ is again an eigenvalue of $F_{v,\rho}$.  Hence, there
exists an index $j$ such that
\[
\psi_j(f_v)=z.x=z.\psi_i(f_v).
\]
But $f_v\in K^*U_m$ and all $\psi_j$ are trivial on $K^*$ and map
$U_m$ into the open subgroup $U$ we started with.  Hence
$z=\psi_j(f_v).\psi_i(f_v)^{-1}$ belongs to $U$, and this
contradicts the way $U$ has been chosen.
\end{proof}

\begin{proof}[Proof of the proposition]
Since the $n_\sigma(i)$ are divisible by $N$, there exist
$\phi_i\in X(T)$ with $\phi_i^N=\chi_i^{\mathrm{alg}}$.  If
$x\in K_\ell^*$, we have:
\[
\phi_i(x^{-1})^N=\chi_i^{\mathrm{alg}}(x^{-1})=\chi_i(x)=\psi_i(x)^N
\]
if $x$ is close enough to $1$.  Hence $\phi_i(x)\psi_i(x)$ is an
$N^{\mathrm{th}}$-root of unity when $x$ is close enough to $1$,
and, by continuity, it is equal to $1$ in a neighbourhood of $1$.
Hence, the restriction of $\rho$ to $K_\ell^*$ is locally equal to
$\phi^{-1}$, where $\phi$ is the (algebraic) representation of $T$
defined by the family $(\phi_1,\ldots,\phi_n)$.  The representation
$\phi$, \emph{a priori} defined over $\bar\Q_\ell$, can be defined
over $\Q_\ell$ (and even over $\Q$); this follows, for instance,
from the fact that the family $(\phi_1,\ldots,\phi_n)$ is
\emph{stable} under the action of $\Gal(\bar\Q/\Q)$, since the
family $(\chi_1^{\mathrm{alg}},\ldots,\chi_n^{\mathrm{alg}})$ is.

Hence $\rho$ is locally algebraic, q.e.d.
\end{proof}

\subsechead{3.3}{An auxiliary result on tori}

In \cite{lang-transc}, Lang proved that two exponential functions
$\exp(b_1z)$, $\exp(b_2z)$, $b_1,b_2\in\C$, which take algebraic
values for at least 3 $\Q$-linearly independent values of $z$, are
multiplicatively dependent: the ratio $b_1/b_2$ is a rational
number.  This had also been noticed by Siegel.

Lang proved the following \textit{l}-adic analogue:

\begin{proposition}[1]\label{prop:III.3.3.1}
%% NOTE(p115): the field is printed "$\Q_\ell$" (no bar) in both places;
%% checked at dpi 500.
Let $E$ be a field containing $\Q_\ell$ and complete for a real
valuation extending the valuation of $\Q_\ell$.  Let $b_1,b_2\in E$
and let $\Gamma$ be an additive subgroup of $E$.  Assume:
\begin{alist}[label=(\arabic*)]
\item[(1.)] $\Gamma$ is of rank at least 3 over $\Z$.
\item The exponential series $\exp(z)=\sum z^n/n!$ converges
absolutely on $b_1\Gamma$ and $b_2\Gamma$.
\item For all $z\in\Gamma$ the elements $\exp(b_1z)$ and
$\exp(b_2z)$ are algebraic over $\Q$.
\end{alist}
Then $b_1$ and $b_2$ are linearly dependent over $\Q$ (i.e.
$b_1/b_2$ belongs to $\Q$ if $b_2\neq 0$).
\end{proposition}

For the proof, see \cite{lang-transc}, Appendix, or
\cite{serre-exponentielles}, \S1.

We will apply this result to tori, taking for $E$ the completion of
$\bar\Q_\ell$.  We need a few definitions first:

a/ Let $T$ be an $n$-dimensional torus over $\Q$, with character
group $X(T)$.  As before, we identify $X(T)$ with the group of
morphisms of $T_{/E}$ into $\mathbf{G}_{m/E}$.  We say that $T$
\emph{is a sum of one-dimensional tori} if there exist
one-dimensional subtori $T_i$ of $T$, $1\leq i\leq n$, such that the
sum map $T_1\times\cdots\times T_n\to T$ is surjective (and hence
has a finite kernel).  An equivalent condition is: $X(T)\tens\Q$
\emph{is a direct sum of one-dimensional subspaces stable by}
$\Gal(\bar\Q/\Q)$.

b/ Let $f$ be a continuous homomorphism of $T(\Q_\ell)$ into
$E^*$.  We say that $f$ is \emph{locally algebraic} if there is a
neighbourhood $U$ of $1$ in the \textit{l}-adic Lie group
$T(\Q_\ell)$, and an element $\phi\in X(T)$ such that
$f(x)=\phi(x)$ for all $x\in U$.  We say that $f$ is \emph{almost
locally algebraic} if there is an integer $N\geq 1$ such that $f^N$
is locally algebraic.

c/ Let $S$ be a finite set of prime numbers, and, for each
$p\in S$, let $W_p$ be an open subgroup of $T(\Q_p)$; denote by $W$
the family $(W_p)_{p\in S}$.

Let $T(\Q)_W$ be the set of elements $x\in T(\Q)$ whose images in
$T(\Q_p)$ belong to $W_p$ for all $p\in S$; this is a subgroup of
$T(\Q)$.  With these notations, we have:

\begin{proposition}[2]\label{prop:III.3.3.2}
Let $f\colon T(\Q_\ell)\to E^*$ be a continuous homomorphism.
Assume:
\begin{alist}
\item There exists a family $W=(W_p)_{p\in S}$ such that $f(x)$ is
algebraic over $\Q$ for all $x\in T(\Q)_W$.
\item $T$ is a sum of one-dimensional tori.
\end{alist}
Then $f$ is almost locally algebraic.
\end{proposition}

\begin{proof}
\begin{alist}[label=\roman*)]
\item We suppose first that $T$ \emph{is one-dimensional}, and we
denote by $\chi$ a generator of $X(T)$.  If $\chi$ is invariant by
$\Gal(\bar\Q/\Q)$, $T$ is isomorphic to $\Gm$ and
$T(\Q)\simeq\Q^*$.  If not, $\Gal(\bar\Q/\Q)$ acts on $X(T)$
\emph{via} a group of order 2, corresponding to some quadratic
extension $F$ of $\Q$; the character $\chi$ defines an isomorphism
of $T(\Q)$ onto the group $F_1^*$ of elements of $F$ of norm 1.  In
both cases, one sees that $T(\Q)$ is an abelian group of
\emph{infinite rank} (for a more precise result, see Exercise
below).  On the other hand, each quotient $T(\Q_p)/W_p$ is a
finitely generated abelian group of rank $\leq 1$.  Hence
$T(\Q)/T(\Q)_W$ is finitely generated, and this implies that
$T(\Q)_W$ is also \emph{of infinite rank}.

Since $T(\Q_\ell)$ is an \textit{l}-adic Lie group of dimension 1,
it is locally isomorphic to the \emph{additive group} $\Q_\ell$.
This means that there exists a homomorphism
\[
e\colon \Z_\ell\to T(\Q_\ell)
\]
which is an \emph{isomorphism} of $\Z_\ell$ onto an open subgroup
of $T(\Q_\ell)$.  By composition we get two continuous
homomorphisms
\[
f\circ e\colon \Z_\ell\to E^*,\qquad
\chi\circ e\colon \Z_\ell\to E^*.
\]
But any continuous homomorphism of $\Z_\ell$ into $E^*$ is locally
an exponential.  This implies that, after replacing $\Z_\ell$ by
$\ell^m\Z_\ell$ if necessary, there exist $b_1,b_2\in E$ such that
\[
f\circ e(z)=\exp(b_1z),\qquad
\chi\circ e(z)=\exp(b_2z),
\]
with absolute convergence of the exponential series.

Let now $\Gamma$ be the set of elements $z\in\Z_\ell$ such that
$e(z)\in T(\Q)_W$.  Since $T(\Q_\ell)/e(\Z_\ell)$ is finitely
generated, and $T(\Q)_W$ is of infinite rank, $\Gamma$ is of
infinite rank.  If $z\in\Gamma$, $e(z)$ belongs to $T(\Q)_W$, hence
$f\circ e(z)$ is algebraic over $\Q$; the same is true for
$\chi\circ e(z)$ since $\chi$ maps $T(\Q)$ either into $\Q^*$ or
into the group $F_1^*$ defined above.
Proposition~\ref{prop:III.3.3.1} then shows that $b_1/b_2$ is
rational.  This means that some integral power $f^N$ of $f$, with
$N\geq 1$, is locally equal to an integral power of $\chi$, hence
$f$ is \emph{almost locally algebraic}.

\item \emph{General case.}  Write $T=T_1\cdots T_n$ where
$T_1,\ldots,T_n$ are one-dimensional subtori of $T$.  Since
$X(T)\tens\Q$ is the direct sum of the $X(T_i)\tens\Q$, it is
enough to show that, for all $i$, the restriction $f_i$ of $f$ to
$T_i(\Q_\ell)$ is almost locally algebraic.  But we may choose open
subgroups $W_{i,p}$ of $T_i(\Q_p)$ such that
$W_{1,p}\cdots W_{n,p}\subset W_p$.  If we put
$W_i=(W_{i,p})_{p\in S}$, we then see that $f_i$ takes algebraic
values on $T_i(\Q)_{W_i}$, hence is almost locally algebraic by i)
above, q.e.d.
\end{alist}
\end{proof}

\runhead{Remark}

If one could suppress condition (b) from
Prop.~\ref{prop:III.3.3.2}, all the results of this \S{} would
extend to arbitrary number fields.  This would be possible if one
had a sufficiently strong $n$-dimensional version of
Prop.~\ref{prop:III.3.3.1} above; the one given in
\cite{serre-exponentielles}, \S2 does not seem strong enough (it
requires density properties which are unknown in the case
considered here).  $\rightarrow$ [This has been done by
Waldschmidt: see \cite{henniart}, \cite{waldschmidt}.]

\runhead{Exercise}

Let $T$ be a non-trivial torus over $\Q$.  Show that $T(\Q)$ is the
direct sum of a finite group and a free abelian group of infinite
rank.

\subsechead{3.4}{Proof of the theorem}

We go back to the notations and hypotheses of \S\ref{sec:III.3.1}.
Let
\[
\rho\colon \Gal(\bar K/K)^{\mathrm{ab}}\to\Aut(V_\ell)
\]
be a rational, semi-simple, abelian \textit{l}-adic representation
of $K$.  If $E$ is the completion of $\bar\Q_\ell$, as in
\S\ref{sec:III.3.3}, we may bring $\rho$ in diagonal form over $E$.
This gives rise to a family $(\psi_1,\ldots,\psi_n)$ of continuous
characters of $\Gal(\bar K/K)^{\mathrm{ab}}$ (hence also of the
idèle group $I$) into $E^*$; here, $n=\dim.\,V_\ell$.

Let $f_i\colon K_\ell^*\to E^*$ be the restriction of $\psi_i$ to
the $\ell^{\mathrm{th}}$-component $K_\ell^*$ of $I$.  Note that
$K_\ell^*=T(\Q_\ell)$, where $T$ is, as usual, the torus defined by
$K$ (chap.~\ref{chap:II}, \S\ref{sec:II.1.1}).

\begin{lemma}
The torus $T$ and the homomorphism $f_i$ satisfy the assumptions
(a) and (b) of Prop.~\ref{prop:III.3.3.2}, \S\ref{sec:III.3.3}.
\end{lemma}

\begin{proof}[Verification of (a)]
Let $S$ be a finite set of primes, with $\ell\nmid S$, such that if
$v\in\Sigma_K$, $p_v\neq\ell$, $p_v\nmid S$, the representation
$\rho$ is unramified at $v$, and the characteristic polynomial of
$F_{v,\rho}$ has rational coefficients.  If $p\in S$, Prop.~1 of
\S\ref{sec:III.2.2} shows that there exists an open subgroup $W_p$
of $K_p^*=T(\Q_p)$ such that $\psi_i(W_p)=1$.  Let
$W=(W_p)_{p\in S}$ and let $x\in T(\Q)_W$.  Since $x\in K^*$, we
have $\psi_i(x)=1$, when $x$ is identified with an idèle of $K$.
On the other hand, let us split the idèle $x$ in its components
\[
x=x_\infty\cdot x_\ell\cdot x_S\cdot x'
\]
according to the decomposition of $I$ in
\[
I=K_\infty^*\times K_\ell^*\times K_S^*\times I'.
\]
(Here $K^*=(K\tens\R)^*$,
$K_S^*=\prod\limits_{p\in S}K_p^*$ and $I'$ is the restricted
product of the $K_v^*$, for $v\in\Sigma_K$, and $p_v\neq\ell$,
$p_v\nmid S$.)  The relation $\psi_i(x)=1$, together with
$\psi_i(x_\ell)=f_i(x)$, gives
\[
f_i(x)^{-1}=\psi_i(x_\infty)\psi_i(x_S)\psi_i(x').
\]
By construction, we have $\psi_i(x_S)=1$ and it is clear that
$\psi_i(x_\infty)=\pm 1$.  Hence:
\[
f_i(x)=\pm\psi_i(x')^{-1}.
\]
But, for each $v\in\Sigma_K$, with $p_v\nmid S$, $p_v\neq\ell$, we
know that the eigenvalues of $F_{v,\rho}$ are algebraic; hence, if
$f_v$ is an idèle which is a uniformizing element at $v$, and is
equal to $1$ elsewhere, $\psi_i(f_v)$ is \emph{algebraic}.  If
$a(v)$ is the valuation of $x'$ at $v$, we have:
\[
\psi_i(x')=\prod\psi_i(f_v)^{a(v)};
\]
hence $\psi_i(x')$ and $f_i(x)$ are algebraic and we have checked
(a).
\end{proof}

\begin{proof}[Verification of (b)]
Since $K$ is a composite of quadratic fields, it is a Galois
extension of $\Q$, and its Galois group $G$ is a product of groups
of order 2.  The character group $X(T)$ of $T$ is isomorphic to the
regular representation of $G$, and it is clear that $X(T)\tens\Q$
splits as a direct sum of one-dimensional $G$-stable subspaces
(each correspond to a character of $G$).  Hence $T$ is a sum of
one-dimensional tori.
%% NOTE(p121): "each correspond to a character of G" is AS PRINTED
%% (singular/plural mismatch in the original).
\end{proof}

%% NOTE(p121): the three blocks below are printed at the top of III-30
%% (PDF 121), before the "APPENDIX" head; they conclude §3.4.

\begin{proof}[End of the proof of the theorem]
Using prop.~\ref{prop:III.3.3.2} of \S\ref{sec:III.3.3}, we see
that each $f_i$ is \emph{almost locally algebraic}.  Hence there is
an integer $N\geq 1$ such that the $f_i^N$ are locally algebraic.
This implies, cf.~\S\ref{sec:III.1.1}, that $\rho^N$ is locally
algebraic, hence (cf.~\S\ref{sec:III.3.2}) that $\rho$ itself is
locally algebraic, q.e.d.
\end{proof}

\runhead{Exercise}

Assume that $K$ is a composite of quadratic fields.  Let $\chi$ be
a \emph{Grössencharakter} of $K$ and suppose that the values of
$\chi$ (on the ideals prime to the conductor) are algebraic
numbers.  Show that $\chi$ is ``of type (A)'' in the sense of Weil
\cite{weil-characters}.  (Use the same method than above, with $E$
replaced by $\C$.)  If the values of $\chi$ lie in a finite
extension of $\Q$, show that $\chi$ is ``of type $(A_0)$''.
$\rightarrow$ [no assumption on $K$ is necessary, thanks to
\cite{waldschmidt}.]
%% NOTE(p121): "the same method than above" is AS PRINTED.
